\chapter{Conclusão}
\label{cap:conclusao}

Este trabalho realizou um estudo comparativo experimental de três implementações do
Quicksort --- recursivo puro, híbrido com Insertion Sort e híbrido com
mediana-de-três --- avaliadas sobre cinco massas de teste e quatro tamanhos de
entrada, com medição de tempo de execução, número de comparações e número de trocas.

Quanto ao \textbf{primeiro objetivo específico}, a calibração empírica indicou
$M \approx 30$--$40$ como faixa ótima, com platô de bom desempenho entre 20 e 70. O
valor adotado, $M = 40$, mostrou-se estável entre execuções independentes, entre
tamanhos de entrada e entre estratégias de pivô. Esse valor é superior à faixa de 5 a
25 recomendada por \citeonline{sedgewick1978}, diferença atribuída à evolução do
hardware: em processadores atuais, a localidade de cache e a previsibilidade de
desvios do Insertion Sort deslocam o ponto de equilíbrio para valores maiores.

Quanto ao \textbf{segundo objetivo}, as versões híbridas reduziram o tempo de
execução entre 16\% e 45\% em relação ao Quicksort recursivo puro, com mediana de
31\%, sem prejuízo em nenhuma configuração testada. O ganho foi maior nas massas com
estrutura, em razão da adaptatividade do Insertion Sort.

Quanto ao \textbf{terceiro objetivo}, o experimento de pior caso confirmou
numericamente a previsão teórica: a razão comparações/$n^2$ convergiu para exatamente
$0{,}5000$, correspondendo às $n(n-1)/2$ comparações previstas pela análise. Para
$n = 50.000$, a diferença entre a escolha inadequada e a adequada de pivô alcançou
aproximadamente 2.000 vezes em número de comparações e 1.700 vezes em tempo de
execução.

Quanto ao \textbf{quarto objetivo}, a análise crítica produziu dois resultados que
merecem destaque.

O primeiro é que a mediana-de-três, quando comparada ao pivô central \textbf{com o
mesmo valor de $M$}, não apresenta vantagem sistemática: reduz as comparações em 4\%
a 9\% apenas em dados aleatórios e é cerca de 1\% pior em entradas ordenadas,
inversas e quase ordenadas. O resultado é plenamente explicado pela teoria --- em um
vetor ordenado o elemento central já é a mediana exata do subvetor, de modo que a
técnica seleciona o mesmo pivô e gasta comparações adicionais sem contrapartida. Seu
valor real reside na proteção contra pivôs de posição fixa, cenário em que a
diferença é de três ordens de grandeza. Cabe destacar que esse resultado só se
tornou visível porque as duas versões foram comparadas mantendo $M$ constante;
compará-las com valores distintos de $M$ atribuiria à mediana-de-três um ganho que
decorre, na verdade, do corte.

O segundo é que a versão híbrida executa \textit{mais} comparações e \textit{mais}
trocas que o Quicksort recursivo em dados aleatórios e, ainda assim, é
consistentemente mais rápida. Isso evidencia que a contagem de operações
elementares, embora seja a métrica correta para a análise assintótica, deixou de ser
um preditor confiável do tempo de execução em hardware moderno, no qual efeitos de
cache e de predição de desvios respondem por parcela expressiva do custo real --- em
consonância com \citeonline{kushagra2014} e \citeonline{edelkamp2019}.

Registra-se, por fim, uma contribuição de natureza metodológica. A primeira
calibração realizada produziu um valor incorreto de $M$ em virtude do aquecimento do
processador, viés sistemático que não é corrigido pela simples repetição das
execuções. O erro foi detectado pela discordância entre o tempo medido e o número de
comparações --- métrica determinística e imune ao efeito. Conclui-se que os
contadores exigidos pelo enunciado desempenham papel que extrapola a caracterização
de desempenho: constituem instrumento de verificação da validade da própria medição
experimental.

\section{Trabalhos futuros}

Como desdobramentos naturais deste estudo, sugerem-se: a implementação do
\textit{Introsort} de \citeonline{musser1997}, que garante $O(n\log n)$ no pior caso
mediante monitoramento da profundidade de recursão; a avaliação do Quicksort de pivô
duplo analisado por \citeonline{wild2012}; e a implementação das técnicas de partição
sem desvios propostas por \citeonline{edelkamp2019} e \citeonline{peters2021}, cujo
ganho, segundo os autores, decorre de fatores que a análise assintótica clássica não
captura --- hipótese que os resultados da Seção \ref{sec:analise-critica} deste
trabalho corroboram.
